\uuid{7SuV}
\titre{Densité}
\auteur{bodin ; hanani ; mzari}
\organisation{exo7}
\licence{CC BY-NC-SA 4.0}
\datecreate{2020-01-20}
\niveau{L1}
\module{Analyse}
\chapitre{Continuité, limite et étude de fonctions réelles}
\sousChapitre{Limite de fonctions}
\difficulte{3}
\theme{Limites des fonctions réelles, densité}
\exo7qcmid{lille-1.limites.023}

\contenu{
\begin{question}
\qclassification{123.03}
\qtags{motcle=densité}
Soit  $f$ une fonction définie sur $[0,1]$ par : $f(x)=\left\{\begin{array}{cc}x-1,& \mbox{si} \, x \in \Rr \setminus \Qq\\ 1,&  \mbox{si} \, x \in \Qq  \end{array}\right.$. Quelles sont les assertions vraies ?
\begin{answers}

    \bad{$\lim_{x\to 0} f(x)=1$}

    \bad{$\lim_{x\to 0} f(x)=0$}

    \bad{$\lim_{x\to 0} f(x)=-1$}

    \good{$f$ n'admet pas de limite en $0$.}

\end{answers}
\begin{explanations}
L'ensemble des rationnels est dense dans $\Rr$. Donc il existe une suite de rationnels $(u_n)$ qui tend vers $0$ et donc  $\lim_{n\to +\infty} f(u_n)=1$.  D'autre part, l'ensemble des irrationnels est dense dans $\Rr$.  Donc il existe une suite d'irrationnels $(v_n)$ qui tend vers $0$ et donc  $\lim_{n\to +\infty} f(v_n)=\lim_{n\to +\infty} (v_n-1) = -1 $. On en déduit que  $f$ n'admet pas de limite en $0$.
\end{explanations}
\end{question}
}
